\documentclass[a4paper,11pt]{article}
\usepackage{jheppub} % for details on the use of the package, please see the JINST-author-manual
\usepackage{lineno}
\linenumbers



\arxivnumber{1234.56789} % if you have one

\title{\boldmath A title with some math: $x=1$}

% Collaborations

%% [A] If main author
%% \collaboration{\includegraphics[height=17mm]{collabroation-logo}\\[6pt]
%%  XXX collaboration}

%% or
%% [B] If "on behalf of"
%% \collaboration[c]{on behalf of XXX collaboration}


% Authors
% The "\note" macro will give a warning: "Ignoring empty anchor...", you can safely ignore it.

%% [A] simple case: 2 authors, same institution
%% \author[1]{A. Uthor\note{Corresponding author.}}
%% \author{and A. Nother Author}
%% \affiliation{Institution,\\Address, Country}

%% or, e.g.
%% [B] more complex case: 4 authors, 3 institutions, 2 footnotes
%% \author[a,b]{F. Irst,\note{Now at another university}}
%% \author[c]{S. Econd,}
%% \author[a,2]{T. Hird\note{Also at Some University.}}
%% \author[c,2]{and Fourth}
%% \affiliation[a]{Institution_1,\\Address, Country}
%% \affiliation[b]{Institution_2,\\Address, Country}
%% \affiliation[c]{Institution_3,\\Address, Country}

\author{A. Uthor}
\affiliation{One University,\\
some-street, Country}
\affiliation{Another University,\\
different-address, Country}

% E-mail addresses: only for the corresponding author
\emailAdd{first@one.univ}

\abstract{Abstract...}



\begin{document}
\maketitle
\flushbottom

\section{Constraints from LHC}
\label{sec:intro}

In this section, we summarize the current LHC limits on New Physics particles. 
The results are based on data at a center-of-mass energy of 13 TeV, corresponding to an integrated luminosity of 138-140 fb$^{-1}$. 

\subsection{$Z^\prime$ and $W^\prime$}
In the early 2026, CMS has published their combined statistical analysis of searches for heavy vector boson resonances~\cite{CMS:2026fjw}. 
Both Drell–Yan (DY) and Vector Boson Fusion (VBF) data are analyzed.
For $Z^\prime$, the decay channels $\ell^+\ell^-$, $q\bar{q}$, $t\bar{t}$, $ZZ$, $Zh$ are considered. 
For $W^\prime$, the decay channels $\ell\bar{\nu_\ell}$ ($\ell=e,\mu,\tau$), $q\bar{q}^\prime$, $t\bar{b}$, $WZ$, $Wh$ are considered. 
This combined analysis gives the most stringent constraints of Heavy Vector Triplet (HVT), 
where $V^\prime$ are excluded with masses below $5.5$ TeV.

For sequential standard model, the most stringent limits on $Z^\prime$ bosons come from high-mass dilepton resonance searches, 
which set a lower mass limit on $5.2$ TeV~\cite{CMS:2021ctt}, 
while the $W^\prime$ has a lower limit on $5.7$ TeV, coming from $\ell\bar{\nu_\ell}$ channel~\cite{CMS:2022krd}. 

We summarize these constraints in Table \ref{table:Vprime}. 

\begin{table}
  \centering 
    \begin{tabular}{cccc}
    %{m{15mm} m{70mm} m{18mm}}
    Process  & Decay & Lower Mass Limit & Reference \\
     \hline
    $pp\to W^\prime/Z^\prime$ & $V^\prime\to WH/WZ/ff$ & $5.5$ TeV & \cite{CMS:2026fjw} \\
$pp\to Z^\prime$  & $Z^\prime\to\ell^+\ell^-$ &   
$5.2$ TeV       &     \cite{CMS:2021ctt}        \\%new row
$pp\to W^\prime$  & $W^\prime\to\ell\bar{\nu_\ell}$ &   
$5.7$ TeV       &     \cite{CMS:2022krd}         \\%new row
    \hline
    \end{tabular}
      \caption{Current most stringent experimental limits on $Z^\prime$ and $W^\prime$ from LHC.}
      \label{table:Vprime}
  \end{table} 

\subsection{$H^\pm$}

The searches for $H^\pm$ are categorized by its mass relative to the top quark: light ($m_{H^\pm} < m_t$) and heavy ($m_{H^\pm} > m_t$).

For the light $H^\pm$, top quark can decay to $H^\pm$ and $b$ quark, 
and $H^\pm$ further decays to Standard Model particles. 
ATLAS has reported two decay channels: $H^+ \to \tau^+\nu_\tau$~\cite{ATLAS:2024hya}, $H^\pm \to cs$~\cite{ATLAS:2024oqu} and $H^\pm\to cb$~\cite{ATLAS:2023bzb}. 
For $H^\pm \to \tau\nu_\tau$, current limits on the branching fraction $\mathcal{B}(t \to H^\pm b) \times \mathcal{B}(H^\pm \to \tau\nu_\tau)$ are as low as $0.02\%$ to $0.27\%$ with $m_{H^\pm}$ from $80$ GeV to $160$ GeV.
For $H^\pm \to cs$, upper limits on $\mathcal{B}(t \to H^\pm b)$ in the $cs$ channel between $0.07\%$ and $3.6\%$, 
assuming $\mathcal{B}(H^\pm \to cs)=1$ and $m_{H^\pm}\in (60, 168)$ GeV.
Similar to $H^\pm \to cs$, $H^\pm\to cb$ gives the upper limits on $\mathcal{B}(t \to H^\pm b)$ between $0.15\%$ and $0.42\%$, where $m_{H^\pm}\in (60, 160)$ GeV.

For the heavy $H^\pm$, the $pp\to tbH^\pm$ has been widely studied since it is the dominant process of 2HDM model. 
The decay channels $H^\pm \to tb$~\cite{ATLAS:2021upq}, $H^\pm \to Wh$~\cite{ATLAS:2024rcu}, and $H^\pm \to \tau\nu_\tau$~\cite{ATLAS:2024hya}. 
The mass range from $200$ GeV to $3$ TeV has been searched. 
The constraints for $H^\pm$ are summarized in Table.~\ref{table:chargehiggs}.
\begin{table}
  \centering 
    \begin{tabular}{cccc}
    %{m{15mm} m{70mm} m{18mm}}
    Process  & Decay & Mass Range & Reference \\
     \hline
$t\to H^\pm b$  & $H^+ \to \tau^+\nu_\tau$ &   
$80-160$ GeV       &   \cite{ATLAS:2024hya}        \\%new row
$t\to H^\pm b$  & $H^\pm\to cs$ &   
$60-168$ GeV       &    \cite{ATLAS:2024oqu}        \\%new row
$t\to H^\pm b$  & $H^\pm\to cb$ &   
$60-160$ GeV       &    \cite{ATLAS:2023bzb}        \\%new row
    \hline
    $pp\to tbH^\pm$ & $H^\pm\to tb$ & 200-2000 GeV & \cite{ATLAS:2021upq} \\
    $pp\to tbH^\pm$ & $H^\pm\to Wh$ & 230-3000 GeV & \cite{ATLAS:2024rcu} \\
    $pp\to WbH^\pm b$ & $H^\pm\to \tau\nu_\tau$ & 80-3000 GeV & \cite{ATLAS:2024hya} \\
    \hline
    \end{tabular}
      \caption{Current most stringent experimental limits on $H^\pm$ from LHC.}
      \label{table:chargehiggs}
  \end{table} 

\subsection{$H^{\pm\pm}$}

At LHC, doubly charged Higgs bosons can be produced via DY process or VBF process. 
Then they can decay to same signed leptons or vector bosons. 
Ref.~\cite{ATLAS:2022pbd} reports the search for DY production and the leptonic decay of $H^{\pm\pm}$. 
All combinations of lepton pairs have been considered, and they have the same branching ratio ($1/6$). 
This analysis gives a lower mass limit on $1.08$ TeV.
In another analysis~\cite{ATLAS:2024txt}, the VBF production and $H^{\pm\pm}\to W^\pm W^\pm$, 
where $W$ further decays to lepton and neutrino, 
is studied. 
In this case, the mass range from $0.2$ TeV to $1.5$ TeV is excluded.
The searches for doubly charged Higgs bosons are listed in table.~\ref{table:doublycharge}.

\begin{table}
  \centering 
    \begin{tabular}{cccc}
    %{m{15mm} m{70mm} m{18mm}}
    Process  & Decay & Lower Mass Limit & Reference \\
     \hline
$pp\to H^{++}H^{--}$  & $H^{\pm\pm}\to \ell^\pm\ell^\pm$ &   
$1.08$ TeV       &     \cite{ATLAS:2022pbd}        \\%new row
$pp\to H^{\pm\pm} jj$  & $H^{\pm\pm}\to W^\pm W^\pm$ &   
$1.5$ TeV       &     \cite{ATLAS:2024txt}        \\%new row
    \hline
    \end{tabular}
  \caption{Current most stringent experimental limits ($E_{cm}=13$ TeV, $\mathcal{L}=139$ fb$^{-1}$)}
  \label{table:doublycharge}
  \end{table}    
~
~
~
~

For internal references use label-refs: see section~\ref{sec:intro}.
Bibliographic citations can be done with "cite": refs.~\cite{a,b,c}.
When possible, align equations on the equal sign. The package
\texttt{amsmath} is already loaded. See \eqref{eq:x}.
\begin{equation}
\label{eq:x}
\begin{aligned}
x &= 1 \,,
\qquad
y = 2 \,,
\\
z &= 3 \,.
\end{aligned}
\end{equation}
Also, watch out for the punctuation at the end of the equations.


If you want some equations without the tag (number), please use the available
starred-environments. For example:
\begin{equation*}
x = 1
\end{equation*}

\section{Figures and tables}

All figures and tables should be referenced in the text and should be
placed on the page where they are first cited or in
subsequent pages. Positioning them in the source file
after the paragraph where you first reference them usually yield good
results. See figure~\ref{fig:i} and table~\ref{tab:i} for layout examples. 
Please note that a caption is mandatory  and it must be placed at the bottom of both figures and tables.

\begin{figure}[htbp]
\centering
\includegraphics[width=.4\textwidth]{example-image-a}
\qquad
\includegraphics[width=.4\textwidth]{example-image-b}
\caption{Always give a caption.\label{fig:i}}
\end{figure}

\begin{table}[htbp]
\centering
\begin{tabular}{lr|c}
\hline
x&y&x and y\\
\hline
a & b & a and b\\
1 & 2 & 1 and 2\\
$\alpha$ & $\beta$ & $\alpha$ and $\beta$\\
\hline
\end{tabular}
\caption{We prefer to have top and bottom borders around the tables.\label{tab:i}}
\end{table}

We discourage the use of inline figures (e.g. \texttt{wrapfigure}), as they may be
difficult to position if the page layout changes.

We suggest not to abbreviate: ``section'', ``appendix'', ``figure''
and ``table'', but ``eq.'' and ``ref.'' are welcome. Also, please do
not use \texttt{\textbackslash emph} or \texttt{\textbackslash it} for
latin abbreviaitons: i.e., et al., e.g., vs., etc.


\paragraph{Up to paragraphs.} We find that having more levels usually
reduces the clarity of the article. Also, we strongly discourage the
use of non-numbered sections (e.g.~\texttt{\textbackslash
  subsubsection*}).  Please also consider the use of
``\texttt{\textbackslash texorpdfstring\{\}\{\}}'' to avoid warnings
from the \texttt{hyperref} package when you have math in the section titles.



\appendix
\section{Some title}
Please always give a title also for appendices.





\acknowledgments

This is the most common positions for acknowledgments. A macro is
available to maintain the same layout and spelling of the heading.

\paragraph{Note added.} This is also a good position for notes added
after the paper has been written.


% Bibliography

%% [A] Recommended: using JHEP.bst file
\bibliographystyle{JHEP}
\bibliography{biblio.bib}

%% or
%% [B] Manual formatting (see below)
%% (i) We suggest to always provide author, title and journal data or doi:
%% in short all the informations that clearly identify a document.
%% (ii) please avoid comments such as "For a review'', "For some examples",
%% "and references therein" or move them in the text. In general, please leave only references in the bibliography and move all
%% accessory text in footnotes.
%% (iii) Also, please have only one work for each \bibitem.

% \begin{thebibliography}{99}

% \bibitem{a}
% Author,
% \emph{Title},
% \emph{J. Abbrev.} {\bf vol} (year) pg.

% \bibitem{b}
% Author,
% \emph{Title},
% arxiv:1234.5678.

% \bibitem{c}
% Author,
% \emph{Title},
% Publisher (year).

% \end{thebibliography}
\end{document}
